% META: {"id":"1SPE-GEOREP-CO-027","chapitre":"1SPE-GEOMETRIE-REPEREE","type_objet":"corrige","exercice_ref":"1SPE-GEOREP-EX-027","capacites_codes":["C3"],"sources_inspiration":[],"status":"approved","origin":"generated","status_history":["generated","audited","accepted","approved"],"release_acceptance":"RELEASE_OWNER_BATCH_ACCEPTANCE","acceptance_closure_digest":"sha256:9b3ccf9a81c5520fbb7e03b7b2d3e7bf2057a02834b6d908bf3be5f49f3b553f","acceptance_receipt_digest":"sha256:4fad86f0282e0fa4e0419cca1ad6379ae7af5a280c04a4a90880f90a1c044242"}
% BEGIN-VERIFY
% from sympy import *
% x, y = symbols('x y')
% # Cercle passant par A(0,0), B(4,0), C(0,6)
% # Forme x^2+y^2+Ax+By+C=0. C passe par (0,0): C=0.
% # Par (4,0): 16+4A=0 => A=-4
% # Par (0,6): 36+6B=0 => B=-6
% # Centre: (2,3), rayon = sqrt(4+9) = sqrt(13)
% assert 4**2 + (-4)*4 == 0
% assert 6**2 + (-6)*6 == 0
% assert 4 + 9 == 13
% END-VERIFY
\begin{corrige}{1SPE-GEOREP-EX-027}

\commentaireMarge{On pose $x^2 + y^2 + Ax + By + C = 0$ et on substitue les trois points.}

$A(0 ; 0)$ : $C = 0$.

$B(4 ; 0)$ : $16 + 4A = 0$, soit $A = -4$.

$C(0 ; 6)$ : $36 + 6B = 0$, soit $B = -6$.

L'équation est $x^2 + y^2 - 4x - 6y = 0$.

Centre : $\Omega(2 ; 3)$, rayon $r = \sqrt{4 + 9} = \sqrt{13}$.

\end{corrige}
